% ECONOMICS_PROBLEM: {"id": "G12-221", "title": "Belief aggregation", "jel_code": "G12", "source_id": "OP-0313"}
% !TeX program = xelatex
\documentclass[11pt,a4paper]{article}
\usepackage[margin=23mm]{geometry}
\usepackage{fontspec}
\setmainfont{Latin Modern Roman}
\usepackage{amsmath,amssymb,mathtools}
\usepackage{xcolor,enumitem,booktabs,longtable,array,xurl,hyperref}
\definecolor{muted}{HTML}{53636C}
\hypersetup{colorlinks=true,urlcolor=blue,linkcolor=blue}
\setlength{\parindent}{0pt}
\setlength{\parskip}{5pt}
\newcommand{\R}{\mathbb R}
\newcommand{\E}{\mathbb E}
\newcommand{\N}{\mathbb N}
\newcommand{\ind}{\mathbf 1}
\newcommand{\Var}{\operatorname{Var}}
\newcommand{\Cov}{\operatorname{Cov}}
\newcommand{\supp}{\operatorname{supp}}
\DeclareMathOperator*{\argmin}{arg\,min}
\DeclareMathOperator*{\argmax}{arg\,max}
\newcommand{\entrymeta}[3]{{\sffamily\small\color{muted}\textbf{\texttt{#1}}\quad|\quad #2\par #3\par}}
\newcommand{\Problemref}[1]{\texttt{#1}}
\newcommand{\JELref}[2]{\texttt{#2} (JEL \texttt{#1})}
\begin{document}
\section*{Belief aggregation}
\textbf{Problem ID:} \texttt{G12-221}\quad \textbf{JEL:} \texttt{G12}\par
% BEGIN ECONOMICS STATEMENT
\label{problem:OP-0313}
{\sffamily\small\color{muted}Original subject group: Asset pricing and financial markets\par}
\label{definitions:problem:30007595}
\textbf{Audit classification:} Explicit published open question. Section 6 explicitly requests better belief aggregation and risk-perception evidence. The particular finite-horizon equilibrium, survey kernel and identification problem are editorial.
\subsection*{Definitions and precise research target}
Consider a finite-horizon market with agents \(i=1,\ldots,N\), asset dividends \(D_t\), prices \(P_t\), wealth \(w_{it}\), utility functions and explicitly specified portfolio constraints. Agent \(i\) evaluates uncertainty using a subjective conditional law \(\mathbb P_{it}\), which need not coincide with the objective data-generating law. Survey responses \(b_{it}\) are noisy measurements of these laws under a stated measurement model. An equilibrium price process is generated by optimal demands \(q_{it}\) and clearing \(\sum_iq_{it}=Q_t\), where \(Q_t\) is exogenous supply. For an unconstrained marginal agent, pricing satisfies
\[P_t=E_{\mathbb P_{it}}[M_{i,t+1}(P_{t+1}+D_{t+1})\mid\mathcal F_{it}],\]
with \(M_{i,t+1}\) the agent's marginal-utility discount factor and \(\mathcal F_{it}\) their information. The task is to characterize an empirically estimable aggregation map from the joint distribution of beliefs, risk perceptions, wealth and constraints to prices and volume. Determine conditions under which survey observations identify this map or its relevant derivatives. Explicitly allow the identity of marginal investors to change with prices and constraints. Quantify sensitivity to unobserved nonparticipants, reporting error and selective survey response. Matching average beliefs is insufficient if price-setting agents have different beliefs or risk tolerance.

\textbf{Additional formal definitions and scope (editorial).} Fix dates \(0,\ldots,T\), the last distribution at \(T\), and terminal ex-dividend price \(P_T=0\). Each investor has an increasing concave utility \(u_i\), discount factor \(\beta_i\in(0,1)\), an information filtration \(\mathcal F_{it}\), and a time-consistent subjective probability measure \(\mathbb Q_i\); \(\mathbb P_{it}\) means its regular conditional law rather than an unrestricted sequence of incompatible forecasts. With interior consumption and unconstrained asset choice, \(M_{i,t+1}=\beta_i u_i'(c_{i,t+1})/u_i'(c_{it})\). Portfolio feasibility includes a stated budget, short-sale and borrowing restrictions. Specify a survey kernel \(K_i(db\mid\mathbb Q_i,\mathcal F_{it},w_{it})\) and sampling probabilities, including nonresponse. The price map is from the full cross-sectional distribution of preferences, information, endowments, beliefs and constraints to a selected equilibrium. Trading volume is \(\frac12\sum_i|q_{it}-q_{i,t-1}|\), after distinguishing issuance from secondary trades; holdings alone do not determine volume.

For the identification part, declare an admissible class \(\Theta\) of structural models, including preferences, technologies, shock laws, measurement/sampling rules and any equilibrium selector. If \(O\) denotes the complete observable history and \(T(\theta)\) the target defined above, the identified set is \(\mathcal I_T=\{T(\theta):\theta\in\Theta,\ \mathcal L_\theta(O)=\mathcal L_{\rm obs}(O)\}\); point identification means this set is a singleton. A \(do(a)\) comparison replaces stated policy/technology equations while fixing the declared exogenous law and re-solving the remaining equations. These definitions do not assert that an identification theorem holds without further restrictions.
\subsection*{Variants and assumption changes}
\begin{enumerate}
\item \textbf{Known preferences, unbiased surveys (editorial).} Assume complete participation, observed wealth/constraints and known survey kernel; characterize the equilibrium aggregation map.
\item \textbf{Selective surveys and changing marginal investors (source-motivated extension).} Permit latent risk tolerance and nonresponse, and derive identified price/belief sets rather than averaging surveys directly.
\end{enumerate}
\subsection*{References and source audit}
\begin{enumerate}
\item Author-hosted January 2022 chapter draft verifies authors/title/date; NBER WP 29977 record separately verifies DOI and April 2022 issue. Locator: Section 6, printed pp.35--36 (PDF pp.36--37). Support: Belief aggregation research agenda. Full January 2022 draft accessible. URL: \url{https://bpb-us-w2.wpmucdn.com/voices.uchicago.edu/dist/f/575/files/2022/01/ExpectationsSurvey_13.pdf}.
\end{enumerate}
\begin{enumerate}\raggedright
\item\hypertarget{definitions:source:30007595:1}{} Klaus Adam; Stefan Nagel. 2022. \emph{Expectations Data in Asset Pricing}. Section 6, printed pp.35--36 (PDF pp.36--37).
\url{https://bpb-us-w2.wpmucdn.com/voices.uchicago.edu/dist/f/575/files/2022/01/ExpectationsSurvey_13.pdf}.
DOI: \url{https://doi.org/10.3386/w29977}.
\end{enumerate}

\subsection*{Common conventions: Source mathematical conventions}

These conventions apply to each entry unless explicitly replaced.

\textbf{Models and identification.} A primitive model \(m\) specifies all probability laws, feasible choices, information, equilibrium and measurement rules used by its target. The admissible class \(\mathcal M\) is fixed independently of outcomes. If \(P_O(m)\) is its observable probability law and \(\tau(m)\) the target, its identified set at observable law \(P\) is
\[
 \mathcal I_\tau(P)=\{\tau(m):m\in\mathcal M,\ P_O(m)=P\}.
\]
Point identification means that this set is a singleton. An empty set signals incompatibility of data and assumptions. Coordinate bounds need not describe the joint set. Bounds are sharp only if every retained target is attainable or arbitrarily approachable by compatible models. Finite bounds require explicit restrictions on unobserved quantities; unrestricted confounding, hidden wealth, extrapolation or arbitrary equilibrium selection can yield uninformative sets.

\textbf{Causal interventions.} A potential outcome \(Y_i(p)\) is the outcome under a complete policy regime \(p\), including timing, financing, information and market spillovers. The target population and units are fixed before treatment. With interference, use the complete assignment vector or a justified exposure mapping; individual no-interference notation alone cannot identify an economy-wide intervention. A difference of mean potential outcomes does not require the cross-world joint distribution. Conditioning on post-treatment migration, survival, enrollment or employment changes the estimand and needs separate justification.

\textbf{Statistical statements.} An estimator, confidence set, forecast or test specifies a sampling law, model class and loss/coverage criterion. Uniform \((1-\alpha)\)-coverage means \(\inf_{m\in\mathcal M}\Pr_m\{\tau(m)\in C_n\}\geq1-\alpha\), or an explicitly asymptotic version. The whole parameter space trivially has coverage, so informativeness requires a stated width, risk or power objective. A forecasting improvement is relative to a named benchmark, information set and evaluation population.

\textbf{Equilibrium and optimization.} An equilibrium comprises feasible plans, optimality, beliefs where needed, budget constraints and all market/resource clearing. The primitives or a stated benchmark define each correspondence \(\mathcal E(p;m)\). Nonemptiness is assumed or proved, never inferred just from compact policy sets. A selected-equilibrium target uses a declared selection rule; a robust target quantifies over all permitted equilibria. Use suprema or infima unless attainment is established. An \(\varepsilon\)-optimal policy is feasible and within \(\varepsilon\) of the finite optimum value. Report an empty feasible set separately.

\textbf{Welfare and accounting.} Normative utility, interpersonal normalization, social weights, numeraire, discounting and the use of public receipts are primitives. Behavioral utility need not coincide with the welfare criterion. Household welfare includes ownership income and rents once; adding profits again to compensating variation generally double-counts them. A monetary equivalent is defined by an explicit transfer and price convention, with monotonicity and range conditions. Average effects, mechanism contributions, transfers and welfare losses are different targets. Infinite sums require convergence and dynamic feasible plans require terminal or no-Ponzi conditions.

\textbf{Source versions.} A working-paper date, revision date and journal publication year are distinguished. A DOI for a journal version is not automatically the DOI of an earlier manuscript. Locators refer to the stated version and printed pages where specified. A metadata-only check does not verify a claim about a full-text conclusion. Every entry's source subsection narrows the provenance claim accordingly.
% END ECONOMICS STATEMENT
\end{document}
